% EM structure and teaching revision 20260928
\subsection{从几何求电容的完整流程}
电容不是给定电荷本身的性质，而是导体几何、介质和边界条件共同决定的比例系数。两端电容器依次选定自由电荷 $+Q,-Q$，求介质中的场，再沿从正极到负极的路径算电压，最后消去 $Q$。
例如长为 $\ell$、内外半径 $a,b$ 的同轴圆柱，忽略端部效应，介电常数为 $\varepsilon$，则
\[
 E_r(r)=\frac{Q}{2\pi\varepsilon\ell r},\quad
 U=\int_a^b E_r\dd r=\frac{Q\ln(b/a)}{2\pi\varepsilon\ell},\quad
 C=\frac{2\pi\varepsilon\ell}{\ln(b/a)}.
\]
若间隙 $d=b-a\ll a$，利用 $\ln(1+d/a)\simeq d/a$ 得
$C\simeq\varepsilon(2\pi a\ell)/d$，正好退回展开为平行板后的近似。
多导体问题中，必须先说明哪两部分作为端口、其余导体是接地还是悬浮。一般关系为 $Q_i=\sum_j C_{ij}V_j$；改变悬浮导体的约束，相当于消去一组不同的未知电势，得到的端口等效电容也可能不同。

